\textbf{GRIT's target accuracy improves with (a) a larger smallest clean eigenvalue, degrades with (b) greater noise, and improves with (c) more pairs.} On the Gaussian synthetic task, GRIT gains 32.9 percentage points between the displayed endpoints in (a), loses 80.3 in (b), and gains 4.9 in (c). Accuracy is measured when both cues reverse. GRIT fixes $r=2$; CR strength uses worst accuracy over eight validation groups, including cue reversals, with selection frozen before testing. ERM and true-subspace projection are controls. All panels fix 96 features, 3,000 labels, and ridge $0.1$. Labeled data and target distributions are shared across conditions. Pair differences receive independent isotropic Gaussian noise; $\sigma_\xi$ is its coordinate standard deviation. Panel (a) fixes $k=256$, $\sigma_\xi=0.1$, and $\lambda_{\max}=1$, so clean trace varies. Panels (b,c) fix clean spectrum $(1,1)$; (b) fixes $k=256$ and (c) fixes $\sigma_\xi=0.1$. Lines show means and bands show sample SD over 20 replicates. Zero-noise and nonspanning controls are reported separately. 
